\section{Conclusion}
\label{sec:conclusion}
We introduced SolidSessionBench, an evidence-based benchmark for generating realistic query sequences in decentralized environments. 
By modeling logical sessions, user interests, and iterative refinement, 
it captures client-side performance characteristics that traditional evaluation methods fail to find. 
While using Solid as the primary use case, the underlying concepts apply to approaches such as Linked Web Storage (LWS),
and we will provide an LWS-specific representation once the W3C Working Group finalizes the specifications.

Our evaluation demonstrates that traditional benchmarking methods, such as microbenchmarks and random query sequences, produce misleading estimates of real-world caching performance. 
\todo{Check if true}
By applying SolidSessionBench to caching in LTQP, we uncovered three critical insights. 
First, logical session switching severely degrades cache hit rates, proving that optimal cache sizing and eviction policies must account for session dynamics. 
Second, deep exploratory query refinement introduces significant cache invalidation challenges. 
However, it also exposes unexploited optimization opportunities, such as the 100$\times$ speedups triggered by cost-model unpredictability and cache dynamics when using disk-based caching or combining small caches with cardinality estimation.
Third, cache-based triple pattern cardinality estimation effectively optimizes link traversal queries, as previous executions provide the prior knowledge traditionally absent in LTQP.

This paper provides a clear path forward for optimizing LTQP in decentralized environments. 
Historically, the absence of prior knowledge hindered LTQP query optimization, leading to inefficient execution plans. 
Furthermore, network latency from numerous HTTP requests imposes a strict lower bound on execution times. 
Well-designed caching algorithms can address both bottlenecks simultaneously, offering a promising direction for future optimization literature.

To realize this potential and support realistic user behavior, client-side query engines must move beyond static, session-agnostic caching. 
Future research should prioritize session-aware caching mechanisms that dynamically identify and protect cross-session core entries. 
Additionally, developing intelligent cache warmup strategies and hierarchical architectures, such as combining small 
in-memory quad stores with large disk-based caches, can mitigate the performance penalties of session switching. 
Next, reverse-engineering the extreme speedup anomalies caused by inaccurate cardinality estimates could unlock
highly optimized execution plans tailored specifically for link traversal. 
Finally, future work should explore alternative cache-based cardinality estimation methods,
particularly for evaluating join cardinalities.

To ensure longevity and reuse, SolidSessionBench has been integrated into the core SolidBench ecosystem,
inheriting its established track record of community support and maintenance.
We formalize this in our sustainability plan \url{https://github.com/SolidBench/SolidBench.js/wiki/Sustainability-Plan}.
Due to the \\ seamless integration into SolidBench, its status as the standard community resource 
for evaluating decentralized query execution~\cite{eschauzier2025revisiting,vandenbrande2025incremunica,hanski2025link,tam2024opportunities}
serves as a proxy for the future adoption of this new benchmark when trying to evaluate any client-specific optimization approaches.

The framework is modular, MIT-licensed, and fully documented, with links to all code, data, 
and experiment setups available at this paper-specific hub: \url{https://github.com/RubenEschauzier/simulated-query-sequence-paper}\rt{Could you move all repos you created for this paper that are now under your personal namespace, to the KNoWS GitHub organization?}.  
This \\ extension-friendly approach to development allows researchers to easily modify or extend specific components to meet future requirements
without disrupting the broader evaluation pipeline.
